\section{Limitations}
\label{sec:limitations}

\begin{enumerate}
\item \textbf{Simulation only, small circuits.} Every circuit runs on a statevector
  simulator; nothing here ran on quantum hardware, and noise is modelled by analytic
  surrogates checked against a density-matrix simulation. On the one-dimensional problems
  each circuit is one qubit without entanglement, which a classical computer evaluates
  cheaply, and simulation cost capped the experiments at 6--8 qubits. No claim of quantum
  computational advantage is made or tested.
\item \textbf{Where the method wins.} The eight-copy circuit's strongest result is one
  applied problem. On Poisson a classical model given the same frequencies remains
  $3.3\times$ more accurate, and at high Helmholtz wavenumbers no model learns the
  solution, so parity there says nothing about capability.
\item \textbf{Statistical power and baselines.} Five seeds give a smallest Wilcoxon $p$ of
  0.0625, so ``confirmed'' means all five seeds agree. In the six-problem matrix only the
  frequency-matched model and the random-frequency circuit are size-matched; the MLP and
  Fourier-feature network run at standard sizes. The eight-copy model was selected from a
  small grid of copy counts and learning rates, while classical families received a
  learning-rate search only.
\item \textbf{Design assumptions.} The encoder is restricted to an affine map so that
  $\Omega$ and coverage are well defined; achieved coverage does not by itself predict
  error; and the design card's predicted benefit is a heuristic, not a bound (it
  predicted 2.5\% on heat and zero on Helmholtz $k{=}4,10$, and the circuit beat the MLP
  on none of them).
\end{enumerate}
